% !TEX root = ../main.tex
\section{Limitations and Threats to Validity}
\label{sec:limitations}

The framework has costs and unresolved edges, and a governance argument that does not
state its own is asking for a trust it has not earned. Eight, in descending order of how
much they would change the conclusions.

\paragraph{No empirical validation of the framework's central claims.} This is a
conceptual and architectural paper. The failure shapes it describes are drawn from
deployment experience and are offered as illustrations rather than as measurements; the
prevalence claims are correspondingly weak, and where the text says a failure is common I
mean that I have encountered it repeatedly, not that I have counted it.
Section~\ref{sec:evaluation} specifies the measurements that would settle these
questions, and Appendix~\ref{sec:benchmark} specifies a corpus and case design that would
support them, but neither has been run. In particular: the claim that faithfulness metrics
score misapplied evidence as clean is an analytic consequence of how those metrics are
defined, and has not been demonstrated by running them.

\paragraph{The sensitivity condition is estimated, not established.}
Section~\ref{sec:sensitivity-protocol} sets out a controlled design and is explicit that
it yields a dependence estimate with error in both directions --- false positives from
residual form sensitivity, false negatives from redundant supporting evidence. Grounding
as defined therefore has one conjunct that is, in practice, measured with noise. A
programme reporting a grounding rate should report the protocol and its error properties
alongside, and a rate quoted without them is not comparable across systems.

\paragraph{The defeater search is relative to a search.} By
Section~\ref{sec:why-defeater}, $\NoDefeater$ is bounded to a governed query rather than
quantified over the evidential base. That makes it checkable and it means an undetected
defeater outside the search is not a violation of the definition. Whether the search is
adequate is a corpus-coverage question that the definition cannot answer, and a system
can satisfy the condition while a contrary authority sits in a repository nobody indexed.

\paragraph{Authority ordering is assumed to exist.} The precedence order $\standing$ is
treated as an institutional artefact. Real authority is issue-specific,
jurisdiction-specific, temporally unstable, and sometimes genuinely incomparable; it is
better described as a maintained decision procedure than as a partial order, and an
institution that has not built one cannot evaluate the condition. The framework converts
that absence into a visible missing deliverable, which is progress, but it does not
supply the deliverable.

\paragraph{The framework is silent on how much control is enough.}
Section~\ref{sec:thresholds} derives thresholds by anchoring to the incumbent manual
standard, which presumes such a standard exists and is defensible. Where the manual
process was never held to an evidentiary standard --- which is common --- the anchor is
missing and the framework offers no substitute.

\paragraph{The proportionality argument could license under-control.} Tiering exists
because the full control set is not economic for all traffic. The same argument can be
used to classify consequential work into a lower tier to avoid the cost, and nothing in
this paper prevents that; tier assignment is exactly where a governance regime is
gamed. The dimensions in Section~\ref{sec:tiers} constrain the argument without settling
it.

\paragraph{Applicability may be less crisp than the formalism suggests.} $c$ is treated
as resolvable into jurisdiction, valid time, material facts, and intended use. Real
decisions include cases where the governing regime is itself contested, where the
material facts are what is under determination, and where the intended use changes after
the assertion is made. The formalism handles unresolved context by declining to answer,
which is correct and is also a cost that a system facing many such cases will pay
heavily.

\paragraph{Generalisation beyond the systems in scope is unaddressed.}
Section~\ref{sec:scope} restricts the argument to autoregressive transformer language
models emitting natural-language assertion. Section~\ref{sec:determinism} argues the
structural properties survive the developments visible now, and that argument is about
current architectures rather than a claim about what is possible. An architecture that
maintained content-level representations, closed its asserted set, and gated on evidence
would satisfy the invariants of Section~\ref{sec:invariants}, and the framework would
then be measuring something that holds rather than something that fails.
